% Conclusion and Future Work
% Generated by report/md2tex.py from Report.md -- edit Report.md, then re-run.

\section{Conclusion and Future Work}

Three claims are supported. Substantively, the Hong Kong release network decomposes into cinema blocs that an unsupervised algorithm recovers from co-appearance alone, at modularity 0.75 and with cross-seed stability 0.89, and the local Cantonese-language bloc is identifiable by the people the industry's own histories name as its core. Methodologically, three parallel strategies applied to three workload shapes on one machine produce three different ceilings \textemdash{} 6.93$\times$, 1.67$\times$, 4.34$\times$ \textemdash{} and the differences are explained by dependency structure and memory hierarchy before they are measured. Practically, a workstation-scale implementation with process-level parallelism and cache-aware blocking captures most of the available gain for the scan and counting tiers, and should not be applied to the iterative tier at this graph size.

Future work, in the order we would attempt it: replace Apriori's candidate generation with FP-Growth and re-measure E4, which should move the ceiling from the irregular phase to the counting phase; obtain a country-of-origin label from a source that has one \textemdash{} Wikidata was unreachable from the test environment, but a bulk dump would remove the largest threat to validity; extend the temporal analysis from fixed-label profiles to a genuinely dynamic community method; and test whether bloc position predicts commercial outcomes once box-office data is joined.
